\scanpage{91}
\mspage{84}
% archive sheet number: 90
\[
X_{i_J}\xleftarrow{\unsure{\sim|\simeq}}\varinjlim_{i\in J}X_i
\]
De plus, si $X\in\mathrm{Filt}^{\pi'}$ avec $\pi'>\pi$, on peut prendre $I$ de cardinal $\leq\pi'^\pi$ (donc de cardinal $\leq\pi'$ si $\pi'^\pi=\pi'$, p.ex. si $\pi'=2^{c'}$, avec $c'\geq\pi$).

\uline{Définition} On appelle un \deleted{\unsure{représentation}} \inserted{système} inductif $(X_i)_{i\in I}$ satisfaisant les conditions a), b), c) ci-dessus, un \uline{système inductif bien adapté} (à la filtration cardinale donnée, pour les \unsure{cardinaux|cardinalités} \deleted{\unsure{$\pi_0$}} $\pi_0,\pi$ \deleted{\unsure{fixés|fixées}} \unsure{donnés|données}), et la représentation correspondante
\[
X\xleftarrow{\sim}\varinjlim X_i
\]
une \uline{représentation bien adaptée} de $X$ en termes d'objets de $\mathrm{Filt}^{\pi}$.

\uline{NB} Quand $E$ est stable par petites $\unsure{\varinjlim|\underline{\lim}}$ filtrantes, la situation se simplifie considérablement. On peut alors prendre
\[
\widetilde I=\{I'\subset I\mid\operatorname{Card}I'\leq\pi,\ I'\text{ filtrant}\}
\]
(sans \unsure{s'astreindre à|s'encombrer de} la condition $I'$ grand devant $\pi_0$) et définir encore le système inductif $(X_{I'})_{I'\in\widetilde I}$ par
\[
X_{I'}=\unsure{\lim|\varinjlim}_{i\in I'}X_i.
\]
On a \unsure{encore|alors}
